\section{Nuisance parameters in a profile likelihood}\label{app:D.47}

An exclusion limit computed with every truncation at its central value moves when a truncation moves, and the profile likelihood of \cref{eq:8.8-like} penalizes that motion through its constraint terms. We derive the three statements \cref{sec:21.4-like} uses: how far the limit moves per unit of each nuisance parameter, \cref{eq:21.4-beta}; the theory band that follows when the nuisance parameters are independent and standard normal, \cref{eq:21.4-band}; and the cost of profiling them in the Neyman construction of \cref{sec:21.4-like}, of second order in their sizes, \cref{eq:21.4-profile}.

\subsection{The limit at fixed nuisance parameters}

In bin $j$ the expected count is $N_j(\sigmae,\vartheta) = \sigmae s_j(\vartheta) + N_{\mathrm{bkg},j}$, with $s_j = N_{\mathrm{sig},j}/\sigmae$ the signal per unit cross-section of \cref{eq:21.4-Nsig}. The theory nuisance parameters act only on $s_j$. Up to terms independent of the parameters, the Poisson part of $-2\ln\mathsf{L}$ is $\sum_j\ell_j(N_j)$ with
\begin{equation*}
  \ell_j(N) = 2\bigl(N - N_{\mathrm{obs},j}\ln N\bigr) , \qquad \ell_j'(N) = 2\Bigl(1 - \frac{N_{\mathrm{obs},j}}{N}\Bigr) .
\end{equation*}
We consider data that show no excess, so that the best fit lies at $\hat{\bar\sigma}_e = 0$, the case of an exclusion. Then the best fit does not depend on the nuisance parameters, which multiply a signal that vanishes there, and at fixed $\vartheta$ the one-sided statistic of \cref{app:C.14} is
\begin{equation}
  \mathsf{t}_\vartheta(\sigmae) = \sum_j\Bigl[\ell_j\bigl(\sigmae s_j(\vartheta) + N_{\mathrm{bkg},j}\bigr) - \ell_j(N_{\mathrm{bkg},j})\Bigr] .
  \label{eq:D.47-tfix}
\end{equation}
The limit $\sigmae^{\mathrm{up}}(\vartheta)$ is where $\mathsf{t}_\vartheta$ reaches the threshold $\mathsf{t}^*$ of the construction. Differentiating $\mathsf{t}_\vartheta(\sigmae^{\mathrm{up}}(\vartheta)) = \mathsf{t}^*$ implicitly with respect to $\vartheta_k$ at fixed threshold gives $\partial\ln\sigmae^{\mathrm{up}}/\partial\vartheta_k$ as minus the average of $\partial\ln s_j/\partial\vartheta_k$ over the bins, weighted with $\ell_j'\,\sigmae s_j$.
In the power form \cref{eq:8.8-M}, the bin signal is an integral of the spectrum times $\prod_k(1 + \delta_k)^{\vartheta_k}$, so at $\vartheta = 0$ the logarithmic derivative $\partial\ln s_j/\partial\vartheta_k$ is the average of $\ln(1 + \delta_k)$ over the signal in the bin, which we write $\ln(1 + \bar\delta_{k,j})$. For a normalization template it is exact; for a tilt it holds at first order, and a finite step is used in practice. Hence \cref{eq:21.4-beta}, with the weights
\begin{equation}
  a_j = \frac{\ell_j'\,N_{\mathrm{sig},j}}{\sum_i\ell_i'\,N_{\mathrm{sig},i}}\bigg|_{\sigmae = \sigmae^{\mathrm{up}}(0)} , \qquad \sum_ja_j = 1 .
  \label{eq:D.47-weights}
\end{equation}
A bin with no events has $\ell_j' = 2$, so for zero counts $a_j$ is the share of bin $j$ in the signal. A bin with fewer events than background, $N_{\mathrm{obs},j} \le N_{\mathrm{bkg},j}$, has $\ell_j' \ge 0$ at every $\sigmae \ge 0$ and a nonnegative weight. A bin with an excess can carry a negative weight: raising the signal there improves the fit and weakens the limit.

\subsection{The theory band}

To first order in the nuisance parameters, $\ln\sigmae^{\mathrm{up}}(\vartheta) = \ln\sigmae^{\mathrm{up}}(0) + \sum_k\beta_k\vartheta_k$. With independent standard-normal $\vartheta_k$ the right side is Gaussian with mean $\ln\sigmae^{\mathrm{up}}(0)$ and variance $\sigma_\beta^2 = \sum_k\beta_k^2$, which is \cref{eq:21.4-band}. The envelope that sets one parameter at a time to $\pm1$ moves the limit by $\ee^{\pm|\beta_k|}$.

\subsection{The cost of profiling}

With standard-normal constraints the profiled statistic is
\begin{equation*}
  \mathsf{t}_{\mathrm{p}}(\sigmae) = \min_\vartheta\Bigl[\mathsf{t}_\vartheta(\sigmae) + \sum_k\vartheta_k^2\Bigr] ,
\end{equation*}
since the overall best fit, at $\hat{\bar\sigma}_e = 0$ and $\vartheta = 0$, contributes zero. Let $x = \ln\sigmae$ and $Q(x) = \mathsf{t}_0(\ee^x)$. To first order in the templates, $\partial\mathsf{t}_\vartheta/\partial\vartheta_k = -\beta_kQ'(x)$ by the result above, so that $\mathsf{t}_\vartheta = Q(x - \sum_k\beta_k\vartheta_k)$, with $\beta$ evaluated at $x$. For zero counts or one bin the weights $a_j$, and hence $\beta$, do not depend on $x$; otherwise their dependence on $x$ enters the result at $O(\sigma_\beta^4)$. Only the combination $y = \sum_k\beta_k\vartheta_k$ enters. At fixed $y$ the smallest $\sum_k\vartheta_k^2$ is $y^2/\sigma_\beta^2$, reached for $\vartheta$ parallel to $\beta$, and
\begin{equation*}
  \mathsf{t}_{\mathrm{p}}(\sigmae) = \min_y\Bigl[Q(x - y) + \frac{y^2}{\sigma_\beta^2}\Bigr] .
\end{equation*}
The minimum lies at $Q'(x - y) = 2y/\sigma_\beta^2$, the conditional best fit
\begin{equation}
  y = \tfrac12\sigma_\beta^2\,Q'(x) + O(\sigma_\beta^4) ,
  \label{eq:D.47-yhat}
\end{equation}
a shift of the signal toward smaller values, since the data show no excess.

The Neyman construction of \cref{sec:21.4-like} evaluates the probability of the outcomes at least as extreme as the observed one with the nuisance parameters at this best fit. The expected counts are then those of the nominal model at the shifted cross-section $\ee^{x - y}$, and the probability equals $1 - \mathrm{CL}$ where the shifted cross-section is the nominal limit, $x - y = x_0$. With \cref{eq:D.47-yhat} at $x_{\mathrm{p}} = x_0 + O(\sigma_\beta^2)$,
\begin{equation*}
  x_{\mathrm{p}} - x_0 = \tfrac12\sigma_\beta^2\,Q'(x_0) + O(\sigma_\beta^4) ,
\end{equation*}
which is \cref{eq:21.4-profile} with $\mathsf{t}' = Q'(x_0) = \sum_j\ell_j'N_{\mathrm{sig},j}$ at the nominal limit. For one bin the result is exact. There the best fit obeys $\vartheta_k = \tfrac12\beta_k\ell'(N)\,\sigmae s(\vartheta)$ exactly, and at the limit the counts $N$ and the signal $\sigmae s(\vartheta)$ are those of the classical limit, so $y = \tfrac12\sigma_\beta^2\mathsf{t}'$ with no higher orders. With no events $\mathsf{t}' = 2N^{\mathrm{up}}$ for any background. A known background lowers $\mathsf{t}'$: with no event and $N_{\mathrm{bkg}} = 1$ the classical limit is $N^{\mathrm{up}} = 1.303$ and $\mathsf{t}' = 2.61$, against 4.61 without background. One event on a background of one gives $\mathsf{t}' = 4.29$ and, at $\sigma_\beta = 0.393$, a factor 1.393.

\subsection{Removing cells from the signal}

Removing cells of the $(\vect{q},\omega)$ plane from the signal lowers each $s_j$ by a nonnegative amount. From \cref{eq:D.47-tfix}, $\partial\mathsf{t}_\vartheta/\partial s_j = \sigmae\,\ell_j'(N_j)$, and for $N_{\mathrm{obs},j} \le N_{\mathrm{bkg},j}$ this is nonnegative at every $\sigmae \ge 0$, since $N_j \ge N_{\mathrm{bkg},j}$. Lowering any $s_j$ of such a bin therefore lowers $\mathsf{t}_\vartheta(\sigmae)$ at every $\sigmae$, and the point where it reaches $\mathsf{t}^*$ moves up. The same holds for $\mathsf{t}_{\mathrm{p}}$, a minimum of functions that each decrease, and for the classical limit of one bin, $\sigmae^{\mathrm{up}} = N^{\mathrm{up}}/s$. Removing signal therefore never strengthens a limit set from bins without an excess. In a bin with $N_{\mathrm{obs},j} > N_{\mathrm{bkg},j}$, $\ell_j'$ is negative for $\sigmae s_j < N_{\mathrm{obs},j} - N_{\mathrm{bkg},j}$, and the statement can fail there.
